% =====================================================================
%  ch4_bmodel.tex — §4 B 模型：由组合数据构造的广义 Tate 曲线
% =====================================================================
\section{B 模型：由组合数据构造的广义 Tate 曲线}\label{s:Bmodel}

\subsection{用形式概形表述的广义 Tate 曲线}\label{ss:formalscheme}

给定定义~\ref{def:combidata} 那样的组合数据，下面的构造是 Mumford 构造的一个特例（实际上是它在万有幂级数环上的版本）。先取 $\Z$-概形 $X^0$，它由诸 $X^0_{v}$ 之并得到，其中对 $v\neq v'$ 把 $x_{e/v}$ 与 $x_{e/v'}$ 等同。所得的节点记为 $x_{e}\in X^0$。

选取如下仿射开子集 $U^0_{e},W^0_{v}\subset X^0$。对 $v\in V$，子集 $W^0_{v}$ 是 $X^0_{v}$ 去掉节点。对 $e\in E$，取 $v,v'$ 为 $e$ 的两个端点，并定义 $U^0_{e}$ 为 $X^0_{v}\cup X^0_{v'}$ 去掉除 $x_e$ 之外的所有节点。有同构
\begin{align*}
W_{v}^0&\cong\Spec\Z[t^{\pm 1},(1-t)^{-1}],\\
U_{e}^0&\cong \Spec \Z[t_{e/v},(1-t_{e/v})^{-1},t_{e/v'},(1-t_{e/v'})^{-1}]/(t_{e/v}t_{e/v'}).
\end{align*}
第一个同构当然依赖于 $X_{v}^0$ 上坐标 $t$ 的选取，该坐标在三个标记点处取值 $0,1,\infty$。

现在定义 $\Z[[q_{e},e\in E]]$ 上的形式概形 $\mX$。它模去所有 $q_e$ 后的约化恰为 $X^0$。先取仿射形式概形 $\cU_{e},\cW_{v}$：
\begin{align*}
\cW_{v}&:=\Spf \cO(W_{v}^0)[[q_{f},f\in E]],\\
\cU_{e}&:=\Spf\Z[T_{e/v},(1-T_{e/v})^{-1},T_{e/v'},(1-T_{e/v'})^{-1}][[q_{f},f\in E]]/(T_{e/v}T_{e/v'}-q_{e}).
\end{align*}
容易看出，对 $e/v$、$e/v'$ 有自然同构
\[
\widehat{\cO(\cU_{e})[T_{e/v}^{-1}]}\xrightarrow{\ \sim\ }\cO(\cW_{v}),\qquad
T_{e/v}\mapsto t_{e/v},\quad T_{e/v'}\mapsto \frac{q_{e}}{t_{e/v}}.
\]

这使我们能按显然的方式把所有 $\cU_{e}$ 粘起来，从而得到形式概形 $\mX$。由 Grothendieck 代数化定理容易看出：存在唯一的（在典范同构意义下）$\Z[[\{q_{e}\}]]$ 上的代数曲线 $X$，使得 $X$ 模 $q_e$ 的约化与 $X^0$ 等同，且 $X^0$ 在 $X$ 中的形式邻域与 $\mX$ 等同。

然而代数化本质上无法显式写出，我们也不需要它，因为凝聚层范畴与完美复形范畴都可以自然地从形式概形得到，即
\[
\Coh(X)\simeq\Coh(\mX),\qquad \Perf(X)\simeq \Perf(\mX).
\]

\begin{remark}\label{remark:inverting_q}
尽管一般的去心形式概形（诸如 $\mX-X^0$ 这样的对象）不易处理，这里它们与通常概形差别不大。也就是说，若要反演某族 $q_{e_1},\dots,q_{e_l}$（例如所有 $q_{e}$），只需取一个底拓扑空间相同的有环空间 $\mX'$，在仿射子集上定义截面为
\[
\cO_{\mX'}(\cU)=\widehat{\cO_{\mX}(\cU)[(q_{e_1}\dots q_{e_l})^{-1}]}.
\]
于是 $\Coh(\mX')=\Coh(\mX)/(q_{e_1}\dots q_{e_l}\text{-挠})$，$\Perf(\mX')$ 类似。
\end{remark}

以下记 $K$ 为 Novikov 域 $\widehat{\mk[T^{\R}]}$，其中 $\mk$ 是某个系数域。如上，记 $A_e\in\R_{>0}$ 为 A 侧 $2$-球的辛面积。取连续同态
\[
\Z[[\{q_e, e\in E\}]]\to K,\qquad q_e\mapsto T^{A_e}
\]
（若允许 A 模型的体形变，也可取赋值为 $A_e$ 的其他元素），就得到 $X$ 的标量扩张 $X_K$。B 侧将是曲线 $X_K$。

\subsection{Schottky 群胚}\label{ss:schottky}

现在用刚性解析几何来描述曲线 $X_K$。为避免混淆，记
\[
Y_v:=X^0_v\times_{\Z} K\cong\PP^1_K,
\]
并沿用记号 $t_{e/v}$ 表示所选射影坐标。

记 $\pi_1(G)$ 为图 $G$ 的基本群胚。定义函子 $g:\pi_1(G)\to \Sch/K$：它把每个 $v\in V$ 送到 $Y_v$，并把每条连接 $v$ 与 $v'$ 的边 $e$（看作态射 $e:v\to v'$）送到映射 $g_{e/v}:X^0_v\to X^0_{v'}$，后者由
\[
t_{e/v'}\bigl(g_{e/v}(x)\bigr)=\frac{q_e}{t_{e/v}(x)}
\]
给出。

固定顶点 $v_0\in V$，得到\emph{Schottky 群} $\Gamma_{G,v_0}:=\pi_1(G,v_0)$，它忠实地作用在 $Y_{v_0}$ 上。群 $\Gamma_{G,v_0}$ 在 $g=g(X_K)$ 个生成元上自由，其非单位元都通过 $Y_v$ 的双曲变换作用。

若把每个 $Y_v$ 看作刚性解析空间，则定义 $F_v\subset Y_v$ 为 $\pi_1(G,v)$-作用的极限点集（$F_v$ 自然是 Cantor 集；例如见 \cite[Chapter I]{Gerritzen-VdPut}）。于是曲线 $X_K$ 作为刚性解析空间，与集合 $\{Y_v-F_v\}_{v\in V}$ 对 $\pi_1(G)$ 的商等同（这也等于对任一 $v\in V$ 的商 $(Y_v-F_v)/\Gamma_{G,v}$）。

对每个顶点 $v\in V$ 与任意实数 $1>s_{e/v}>|q_e|$（对每条半边 $e/v$），定义开解析子集 $U_{v,\{t_{e/v}\}}\subset X_K$ 为
\[
\{1\geq |t_{e/v}|\geq s_{e/v}\ \text{对 }e/v\}\subset Y_v-F_v
\]
的像。又对任一半边 $e/v$ 与任意 $1>s_1\geq s_2>|q_e|$，定义 $U_{e/v,s_1,s_2}$ 为
\[
\{s_1\geq |t_{e/v}|\geq s_2\}\subset Y_v-F_v
\]
的像。显然，若边 $e$ 连接 $v$ 与 $v'$，则 $U_{e/v,s_1,s_2}=U_{e/v',\frac{|q_e|}{s_2},\frac{|q_e|}{s_1}}$。对两个不同的 $v,v'\in V$ 与族 $\{s_{e/v}\},\{s_{e/v'}\}$，有
\[
U_{v,\{s_{e/v}\}}\cap U_{v',\{s_{e/v'}\}}=\bigsqcup_{e}U_{e/v,\frac{|q_e|}{s_{e/v'}},s_{e/v}},
\]
其中并对连接 $v$ 与 $v'$ 的边进行；若 $s_1<s_2$，约定 $U_{e/v,s_1,s_2}=\emptyset$。

我们主要使用如下开仿射子集：
\[
U_v:=U_{v,\{|q_e|^{\frac34}\}_{e/v}},\qquad
U_e:=U_{e/v,|q_e|^{\frac14},|q_e|^{\frac34}}.
\]
% [Tutor Note: 直觉] 这里的 $U_v$ 与 $U_e$ 是 B 侧（$X_K$ 上）的基本开集：$U_v$ 是顶点 $v$ 附近、把三条「管状」邻域都截到 $|q_e|^{3/4}$ 的部分；$U_e$ 是边 $e$ 上真正横跨颈部的环带。选择 $1/4<3/4$ 这两个指数，是为了让 $U_v$ 与 $U_{v'}$ 的交恰好落在 $U_e$ 里（见上面的交集公式），从而可以写下 Čech 型局部化。依据：原文 §4.2 与本工程 03_术语表。
